\documentclass[11pt,english]{article}
\bibliographystyle{jf}
\usepackage[longnamesfirst]{natbib}
%\usepackage[T1]{fontenc}
\usepackage{fontenc}
\usepackage[utf8]{inputenc}
\usepackage[bottom]{footmisc}
\usepackage{tabularx}
\usepackage{booktabs}
\usepackage{amsmath,enumerate}
\usepackage{amsthm}
\usepackage{dsfont}
\usepackage{amssymb}
\usepackage{mathtools}
\usepackage{mathrsfs}
\usepackage{graphicx}
\usepackage{natbib}
\usepackage[labelfont=bf,justification=justified,font=small,skip=2pt,labelsep=period]{caption}
%\usepackage[labelfont=bf,justification=justified,labelsep=period]{caption}
\usepackage{subcaption}
\usepackage{setspace}
\usepackage{arydshln}
\usepackage{chngpage}
\usepackage{float,appendix}
\usepackage{afterpage}
\usepackage{verbatim}
\usepackage{indentfirst}
\usepackage[margin=1in]{geometry}
\usepackage[usenames, dvipsnames]{color}
\usepackage{hyperref}
\usepackage{adjustbox}
% \usepackage{longtable}
\usepackage{booktabs, makecell, longtable,ltxtable}
\usepackage{lipsum}
\usepackage{pdflscape}
\usepackage{siunitx}
\usepackage{chngcntr}
\usepackage{lipsum}
\usepackage{tikz}
\usepackage{pdflscape}
\usepackage[normalem]{ulem}
\usepackage{chngcntr,apptools}
\AtAppendix{\counterwithin{lemma}{section}}
\usepackage{siunitx}
\usepackage{tabularx}
\sisetup{table-format=-1.3, table-space-text-post={**}}
\usepackage{xr}
\usepackage{etoolbox}
\usetikzlibrary{positioning}
\newcommand\fnsep{\textsuperscript{,}}
\usepackage{tikz}
\usepackage{pdflscape}
\usepackage[normalem]{ulem}
\usepackage{chngcntr,apptools}
\AtAppendix{\counterwithin{lemma}{section}}
\usepackage{siunitx}
\usepackage{tabularx}
\usepackage{xr}
\usepackage{etoolbox}
\usetikzlibrary{positioning}

%% UCLA colors: 
% https://brand.ucla.edu/wp-content/uploads/ucla-color-palette-v10.pdf
%\definecolor{UCLA}{RGB}{50, 132, 191}   % UCLA Blue primary
%\definecolor{UCLA}{RGB}{52, 123, 173}   % UCLA Blue text only
\definecolor{UCLAblue}{RGB}{30, 75, 135} % secondary color	
\definecolor{UCLAgold}{RGB}{255, 232, 0} % UCLA Gold

\definecolor{usccardinal}{rgb}{0.6, 0.0, 0.0}
\definecolor{Matlabblue}{rgb}{0,0.447,0.741}
\definecolor{bleudefrance}{rgb}{0.19,0.55, 0.91}
\definecolor{cobalt}{rgb}{0.0, 0.28,0.67}
\definecolor{darkblue}{rgb}{0.0, 0.0,0.55}
\definecolor{darkcerulean}{rgb}{0.03,0.27, 0.49}
\definecolor{darkpowderblue}{rgb}{0.0,0.2, 0.6}
\definecolor{bleudefrance}{rgb}{0.19,0.55, 0.91}
%\usepackage[colorlinks=true,urlcolor=blue,linkcolor=blue,citecolor=blue]{hyperref}

\hypersetup{
	colorlinks,
	linkcolor={UCLAblue},
	citecolor={UCLAblue},
	urlcolor={UCLAblue},
}

\bibpunct{(}{)}{;}{a}{,}{,}
\newcommand{\E}{\mathbb{E}}
\newcommand{\var}{\mathrm{Var}}
\newcommand{\cov}{\mathrm{Cov}}
\setcounter{MaxMatrixCols}{20}
\newtheorem{theorem}{Theorem}
\newtheorem{lemma}{Lemma}
\newtheorem{defn}{Definition}
\newtheorem{prop}{Proposition}
\newtheorem{cor}{Corollary}
\urlstyle{same}
\usepackage{pgfplots} 
\pgfplotsset{compat=newest} 
\pgfplotsset{plot coordinates/math parser=false} 
% puts figures and tables at the end of the document
%\usepackage[nomarkers,nolists]{endfloat}               
% puts ONLY figures at the end of the document
%\usepackage[nomarkers,nolists,figuresonly]{endfloat}  
% allows more than one figure per page for endfloat
%\renewcommand{\efloatseparator}{\mbox{}}              
\usepackage{rotating}
% allows sideways tables and figures to be placed at the end of the document
%\DeclareDelayedFloatFlavor{sidewaystable}{table}
%\DeclareDelayedFloatFlavor{sidewaysfigure}{figure}

% \usepackage[nomarkers,nolists]{endfloat} 
% \renewcommand{\efloatseparator}{\mbox{}}              % allows more than one figure per page for endfloat

%\usepackage{mathpazo}
%\usepackage{mathpple}
%%%%%%%%%%%%%%%%%%%%%%%%%%%%%%%%
%\usepackage{palatino}
%%%%%%%%%%%%%%%%%%%%%%%%%%%%%%%%
\pgfplotsset{compat=newest} 
\pgfplotsset{plot coordinates/math parser=false} 
\newlength\figureheight 
\newlength\figurewidth 
\usepackage{mathtools}
\usetikzlibrary{calc}


%%%%%%%%%%%%%%%%%%%%%%%%%%
% Notes customization %
%%%%%%%%%%%%%%%%%%%%%%%%%%
\definecolor{webblue}{rgb}{0,0,0.6}
\definecolor{webgreen}{rgb}{0,.6,0}
\definecolor{webbrown}{rgb}{.6,0,0}


%- Commenting
\let\oldmarginpar\marginpar
\renewcommand\marginpar[1]{\-\oldmarginpar[\raggedleft\tiny #1]%
{\raggedright\tiny #1}}
\newcommand{\mknote}[1]{\marginpar{\textcolor{webblue}{\tiny{Mahyar: #1}}}}
\newcommand{\dsnote}[1]{\marginpar{\textcolor{webgreen}{\tiny{Dejanir: #1}}}}

\title{\textbf{A Model with Equity and Risky Debt Claims}}
\author{ Dejanir Silva\thanks{Daniels School of Business, Purdue University, e-mail: \href{dejanir@purdue.edu}{dejanir@purdue.edu}.} \and Mahyar Kargar\thanks{University of Illinois at Urbana-Champaign, e-mail: \href{mailto:kargar@illinois.edu}{kargar@illinois.edu}.}}
% \date{October 2019}
%\date{\today\\
%	\begin{small}
%		First draft: May 5, 2018
%\end{small}}
\begin{document}
	\sloppy
	\maketitle
	\thispagestyle{empty}
	
\doublespacing

In these notes, we extend the baseline model to allow investors to trade a claim on equity and risky debt. 
We proceed in two steps. 
First, we consider a version where a claim on the aggregate endowment is the only asset in positive net supply. 
In addition to the claim on the aggregate endowment, investors can trade a derivative, so markets are dynamically complete. 
Then, we extend the model to allow for investors to trade an equity claim and a risky debt claim, where the dividends and debt payments add up to the aggregate endowment.


\section{The economy with an aggregate claim}

\subsection{Environment}

\paragraph{Asset structure.}
The aggregate endowment, $Y_t$, follows the process:
\begin{equation}
    \frac{d Y_t}{Y_t} = \mu dt + \sigma d Z_t,
\end{equation}
where $Z_t \in \mathbb{R}^d$ is a vector of independent standard Brownian motions. 
The diffusion term $\sigma$ is a $d$-dimensional row vector. 
In the baseline model, we have $d = 2$, capturing fundamental shocks to cash-flows and portfolio demand shocks to passive investors. 

The price of the aggregate claim is $P_{Y,t}$, and it follows the process:
\begin{equation}
    \frac{d P_{Y,t}}{P_{Y,t}} = \mu_{P_Y,t} dt + \sigma_{P_Y,t} d Z_t.
\end{equation}

The return on a claim on the aggregate endowment follows the process:
\begin{equation}
    d R_{Y,t} = \mu_{R_Y,t} dt + \sigma_{R_Y,t} d Z_t,
\end{equation}
where $\mu_{R_Y,t} \equiv \frac{Y_t}{P_{Y,t}} + \mu_{P_Y,t}$ and $\sigma_{R_Y,t} = \sigma_{P_Y,t}$.

Absence of arbitrage opportunities implies the existence of a stochastic discount factor (SDF), $eta_t$, which follows the process:
\begin{equation}
    \frac{d \eta_t}{\eta_t} = -r_t dt - \lambda_t d Z_t,
\end{equation}
where $r_t$ is the risk-free rate and $\lambda_t$ is the vector of market prices of risk.

The risk premium on the aggregate claim is given by
\begin{equation}
    \pi_{Y,t} \equiv \mu_{R_Y,t} - r_t = \sigma_{R_Y,t} \lambda_t'.
\end{equation}

Investors can also trade a derivative in zero net supply. 
The derivative gives a risk exposure $\sigma_{R_F,t}$ for a premium $\pi_{F,t} \equiv \sigma_{R_F,t} \lambda_t'$.
As long as $\sigma_{R_F,t}$ is not proportional to $\sigma_{R_Y,t}$, markets are dynamically complete.
We focus on the case where $\sigma_{R_F,t}$ is orthogonal to $\sigma_{R_Y,t}$, $\sigma_{R_F,t} \sigma_{R_Y,t}' = 0$, and we adopt the normalization $\|\sigma_{R_F,t}\|  = \|\sigma_{R_Y,t}\|$.

Let $\Sigma_t$ denote the matrix of risk exposures:
\begin{equation}
    \Sigma_t = \begin{pmatrix}
        \sigma_{R_Y,t} \\
		 \sigma_{R_F,t}
    \end{pmatrix},
\end{equation}
and $\pi_t \equiv [\pi_{Y,t}, \pi_{F,t} ]'$.

The no-arbitrage condition can be written as
\begin{equation}
    \pi_t = \Sigma_t \lambda_t' \iff \lambda_t'  = \Sigma_t^{-1} \pi_t.
\end{equation}

\paragraph{Investors' problem.} 

An investor of type $j = 0,\ldots,J$ solves
\begin{equation}\label{eq:inv-problem}
    V_{j,t} = \max_{[C_j,\alpha_j,\theta_j]} 
    \mathbb{E}_t \left[ \int_t^\infty f_j(C_{j,s},V_{j,s})\, ds \right], \qquad s.t.
\end{equation}
% subject to the budget constraint  
\begin{equation}
dW_{j,t} = \big[(r_t + \pi_{Y,t} \alpha_{j,t} + \pi_{F,t} \theta_{j,t}) W_{j,t} - C_{j,t}\big] dt 
    + \Sigma_{j,t} W_{j,t} \, dZ_t,
\end{equation}
where 
\begin{equation}
    \Sigma_{j,t} = \alpha_{j,t} \sigma_{R_Y,t} + \theta_{j,t} \sigma_{R_F,t},
\end{equation}
with $W_{j,t} \geq 0$ and $(\alpha_{j,t}, \theta_{j,t}) \in \Omega_{j,t}$. 

For passive investors ($j=0$),  
$\Omega_{0,t} = \{(\alpha_{0,t}, \theta_{0,t}): \alpha_{0,t} = \overline{\alpha}_{p,t}, \theta_{0,t} = 0\}$, where $\overline{\alpha}_{p,t}$ follows 
\begin{equation}
    d \overline{\alpha}_{p,t} = \theta_p (\overline{\alpha} - \overline{\alpha}_{p,t}) dt + \sigma_{p}  d Z_t.
\end{equation}	
For active investors ($j=1,\ldots,J$),  
$\Omega_{j,t} = \{(\alpha_j, \theta_j): \| \Sigma_{j,t} \| \leq \overline{\sigma}\}$. 

\paragraph{Market clearing.} The  market clearing conditions are
\begin{equation}
	\sum_{j=0}^J \omega_j C_{j,t} = Y_t, 
	\qquad 
	\sum_{j=0}^J \omega_j Q_{j,t} = 1, 
	\qquad 
	\sum_{j=0}^J \omega_j B_{j,t} = 0,
	\qquad
	\sum_{j=0}^J \omega_j W_{j,t} \theta_{j,t} = 0,
\end{equation}
where $Q_{j,t} = \alpha_{j,t} W_{j,t} / P_t$ and $B_{j,t} = (1-\alpha_{j,t}) W_{j,t}$ denote holdings of the risky and riskless assets. 

\subsection{Characterization of the equilibrium}

\paragraph{Value function and consumption-wealth ratio.} With homothetic preferences, investor $j$'s value function can be written as
\begin{equation}\nonumber
    V_{j,t} = \left( \frac{c_{j,t}}{\rho^\psi} \right)^{\frac{1-\gamma_j}{1-\psi}} 
    \frac{W_{j,t}^{1-\gamma_j}}{1-\gamma_j}, \qquad \text{where} \qquad \frac{dc_{j,t}}{c_{j,t}} = \mu_{c_j,t}\, dt + \sigma_{c_j,t}\, dZ_t,
\end{equation}
and $c_{j,t} = c_j(X_t)$ corresponds to the consumption-wealth ratio: $\frac{C_{j,t}}{W_{j,t}} = c_{j,t}$.

From the Hamilton-Jacobi-Bellman equation, the consumption-wealth ratio is
\begin{equation}\label{eq:HJB_complete}
    c_{j,t} = \psi \rho 
    + (1-\psi) \left[ r_t + \pi_{Y,t} \alpha_{j,t} + \pi_{F,t} \theta_{j,t}
        - \frac{\gamma_j}{2} \|\Sigma_{j,t}\|^2 \right] 
    + \xi_{j,t},
\end{equation}
where 
\[
    \xi_{j,t} \equiv \mu_{c_j,t} 
    + (1-\gamma_j) \sigma_{c_j,t} \Sigma_{j,t}' 
    + \frac{\psi - \gamma_j}{1-\psi} \frac{\|\sigma_{c_j,t}\|^2}{2}.
\]

\paragraph{Optimal portfolio share and derivative exposure.} 
The active investor chooses $(\alpha_{j,t},\theta_{j,t})$ to maximize the right-hand side of \eqref{eq:HJB_complete} subject to $\|\Sigma_{j,t}\|\le\overline{\sigma}$. 
Denote the Lagrange multiplier on the leverage constraint by $\nu_{j,t}\ge 0$. The FOC with respect to $\alpha_{j,t}$ is
\begin{equation}\label{eq:FOC_alpha_complete}
    (1-\psi)\Big[\pi_{Y,t} - \gamma_j\,\Sigma_{j,t}\,\sigma_{R_Y,t}'\Big] + (1-\gamma_j)\,\sigma_{c_j,t}\,\sigma_{R_Y,t}' = \nu_{j,t}\,\Sigma_{j,t}\,\sigma_{R_Y,t}'.
\end{equation}
The FOC with respect to $\theta_{j,t}$ is
\begin{equation}\label{eq:FOC_theta_complete}
    (1-\psi)\Big[\pi_{F,t} - \gamma_j\,\Sigma_{j,t}\,\sigma_{R_F,t}'\Big] + (1-\gamma_j)\,\sigma_{c_j,t}\,\sigma_{R_F,t}' = \nu_{j,t}\,\Sigma_{j,t}\,\sigma_{R_F,t}'.
\end{equation}

Using the fact that $\pi_{Y,t} = \sigma_{R_Y,t} \lambda_t'$ and $\pi_{F,t} = \sigma_{R_F,t} \lambda_t'$, we can rewrite the FOC as follows:
\begin{equation}
	\left[ \lambda_t - \lambda_{j,t}  \right] \sigma_{R_Y,t}'  = 0, \qquad 
	\left[ \lambda_t - \lambda_{j,t}  \right] \sigma_{R_F,t}'  = 0,
\end{equation}
where
\begin{equation}
	\lambda_{j,t} \equiv \gamma_j \left[\Sigma_{j,t} + \frac{1 - \gamma_j^{-1}}{1-\psi} \sigma_{c_j,t}\right] + \frac{\nu_{j,t}}{1-\psi} \Sigma_{j,t}.
\end{equation}

The optimality conditions above imply that
\begin{equation}
    \lambda_{j,t} = \lambda_t \iff \lambda_{j,t}' = \Sigma_t^{-1} \pi_t.
\end{equation}

\paragraph{Unconstrained solution.} 
When the leverage constraint is not binding, $\nu_{j,t} = 0$, the optimality conditions above imply that
\begin{equation}
    \Sigma_{j,t} = \frac{\lambda_t}{\gamma_j} + \frac{1 - \gamma_j^{-1}}{\psi-1} \sigma_{c_j,t}.
\end{equation}

We can write $\Sigma_{j,t}$ as
\begin{equation}
    \Sigma_{j,t} = 
	\begin{pmatrix}
		\alpha_{j,t}, 
		\theta_{j,t}
	\end{pmatrix}
	\Sigma_t \iff 
	\begin{pmatrix}
		\alpha_{j,t}, 
		\theta_{j,t}
	\end{pmatrix} = \Sigma_{j,t}\Sigma_t^{-1}.
\end{equation}

Equivalently, we can write the $\alpha_{j,t}$ and $\theta_{j,t}$ as
\begin{equation}
    \alpha_{j,t} = \frac{\pi_{Y,t}}{\gamma_j \|\sigma_{R_Y,t}\|^2} + \varsigma_{Y,j,t},
\end{equation}
where $\varsigma_{Y,j,t} \equiv \frac{1 - \gamma_j^{-1}}{\psi-1} \frac{\sigma_{c_j,t} \sigma_{R_Y,t}'}{\sigma_{R_Y,t} \sigma_{R_Y,t}'}$, 
and
\begin{equation}
    \theta_{j,t} = \frac{\pi_{F,t}}{\gamma_j \|\sigma_{R_F,t}\|^2} + \varsigma_{F,j,t},
\end{equation}
where $\varsigma_{F,j,t} \equiv \frac{1 - \gamma_j^{-1}}{\psi-1} \frac{\sigma_{c_j,t} \sigma_{R_F,t}'}{\sigma_{R_F,t} \sigma_{R_F,t}'}$.

\paragraph{Constrained solution.}
When the leverage constraint is binding, $\nu_{j,t} > 0$, the optimality conditions above imply that
\begin{equation}
	\Sigma_{j,t} = \frac{\lambda_t}{\gamma_j} + \frac{1 - \gamma_j^{-1}}{\psi-1} \sigma_{c_j,t}+ \frac{\nu_{j,t}\gamma_j^{-1}}{\psi-1} \Sigma_{j,t},
\end{equation}

We can then write $\Sigma_{j,t}$ as:	
\begin{equation}
    \Sigma_{j,t} = \frac{\overline{\sigma}}{\|\Sigma_{j,t}^*\|}\,\Sigma_{j,t}^{*}.
\end{equation}

\paragraph{Risk premium and interest rate.} 

Consider the case where the constraint is not binding for all active investors.
Let $\mathcal{J}_t^u$ and $\mathcal{J}_t^c$ denote the sets of unconstrained and constrained active investors at time $t$, with wealth shares $x_{u,t} \equiv \sum_{j\in \mathcal{J}_t^u} x_{j,t}$ and $x_{c,t} \equiv \sum_{j\in \mathcal{J}_t^c} x_{j,t}$. From market clearing for the risky asset, demand from unconstrained investors satisfies
\begin{equation}\label{eq:mkt-clearing}
    \underbrace{\sum_{j=1}^J x_{j,t} \alpha_{j,t} }_{\substack{\text{active demand}}} = \underbrace{1 - x_{0,t} \overline{\alpha}_{p,t} }_{\text{net supply}}.
\end{equation}
% \footnote{
% % This condition is the dynamic analogue of the static risky-asset market clearing condition. 
% % since $\omega_j Q_{j,t} = x_{j,t} \alpha_{j,t}$. 
% When all active investors are unconstrained, this condition collapses to the benchmark in Section~\ref{sec:simple}: the demand from active investors must equal the total supply minus the amount held by passive investors.}  

Combining this condition with the optimal portfolios yields the risk premium:
\begin{equation}\label{eq:sharpe_ratio}
    \pi_{Y,t} = \frac{\gamma_{u,t} \|\sigma_{R_Y,t}\|^2}{x_{a,t}}
    \left[ 1 - x_{0,t}\, \overline{\alpha}_{p,t}
    - x_{a,t}\, \varsigma_{Y,t} \right],
\end{equation}
where $\gamma_{u,t} \equiv [\sum_{j=1}^J \frac{x_{j,t}}{x_{a,t}} \frac{1}{\gamma_j}]^{-1}$ is the average risk aversion and $\varsigma_{Y,t} \equiv \sum_{j=1}^J \frac{x_{j,t}}{x_{a,t}} \varsigma_{Y,j,t}$ the average hedging demand of unconstrained investors.  

The market clearing condition for the derivative is:	
\begin{equation}
    \sum_{j=1}^J x_{j,t} \theta_{j,t} = 0.
\end{equation}
Combining this condition with the optimal portfolios yields the risk premium:
\begin{equation}
    \pi_{F,t} = -\gamma_{u,t} \|\sigma_{R_F,t}\|^2
     \varsigma_{F,t},
\end{equation}
where $\varsigma_{F,t} \equiv \sum_{j=1}^J \frac{x_{j,t}}{x_{a,t}} \varsigma_{F,j,t}$ the average hedging demand of unconstrained investors.  
Goods market clearing requires
\begin{equation}
    \sum_{j=0}^J x_{j,t} c_{j,t} = \frac{1}{p_t}.
\end{equation}
The equilibrium interest rate:
\begin{align}\label{eq:interest_rate}
    r_t &= \rho + \psi^{-1} (\mu_{P,t} + \xi_t )
    + (1-\psi^{-1}) \sum_{j=0}^J x_{j,t} 
        \frac{\gamma_j}{2} \|\Sigma_{j,t}\|^2
    - \pi_{Y,t},
\end{align}
where $\xi_t \equiv \sum_{j=0}^J x_{j,t} \xi_{j,t}$.  

\section{The economy with equity and risky debt claims}

\subsection{Environment}

Aggregate endowment follows
\[
    \frac{dY_t}{Y_t}=\mu\,dt+\sigma\,dZ_t,
\]
where $Z_t$ is a $d$-dimensional Brownian motion and $\sigma\in\mathbb{R}^{1\times d}$.

Dividends are a stochastic share of output, $D_t=s_tY_t$, where the share $s_t\in(0,1)$ follows
\begin{equation}\label{eq:s_process}
    ds_t=\kappa_s(\overline{s}-s_t)\,dt + s_t(1-s_t)\,\sigma_s\,dZ_t,
\end{equation}
with $\kappa_s>0$, $\overline{s}\in(0,1)$, and $\sigma_s\in\mathbb{R}^{1\times d}$. 
In this version of the model, we will have $d = 3$, capturing fundamental shocks to cash-flows, portfolio demand shocks to passive investors, and a dividend share shock.

The risky-debt cash flow is the residual
\[
    B_t=(1-s_t)Y_t,
    \qquad\text{so that}\qquad
    D_t+B_t=Y_t.
\]

As shown above, applying It\^{o}'s lemma implies
\begin{equation}\label{eq:D_dynamics_vector}
    \frac{dD_t}{D_t}=\mu_D(s_t)\,dt+\sigma_D(s_t)\,dZ_t,
\end{equation}
where
\begin{equation}\label{eq:D_mu_sigma_vector}
    \sigma_D(s)=\sigma + (1-s)\sigma_s,
    \qquad
    \mu_D(s)=\mu + \kappa_s\!\left(\frac{\overline{s}}{s}-1\right) + (1-s)\,\sigma\sigma_s',
\end{equation}
and $\sigma\sigma_s'$ denotes the scalar inner product.

Similarly, for risky debt,
\begin{equation}\label{eq:B_dynamics_vector}
    \frac{dB_t}{B_t}=\mu_B(s_t)\,dt+\sigma_B(s_t)\,dZ_t,
\end{equation}
with
\begin{equation}\label{eq:B_mu_sigma_vector}
    \sigma_B(s)=\sigma - s\sigma_s,
    \qquad
    \mu_B(s)=\mu - \kappa_s\!\left(\frac{\overline{s}-s}{1-s}\right) - s\,\sigma\sigma_s'.
\end{equation}

\paragraph{Equity and debt claims.} Let $P_{D,t}$ and $P_{B,t}$ denote the ex-dividend prices of claims to dividend and bond cash flows, respectively. Define instantaneous returns on the two risky claims:
\begin{align}
    dR_{D,t} &\equiv \frac{dP_{D,t}+D_t\,dt}{P_{D,t}}
    = \mu_{R_D,t}\,dt+\sigma_{R_D,t}\,dZ_t, \label{eq:RD_def}\\
    dR_{B,t} &\equiv \frac{dP_{B,t}+B_t\,dt}{P_{B,t}}
    = \mu_{R_B,t}\,dt+\sigma_{R_B,t}\,dZ_t. \label{eq:RB_def}
\end{align}
In a Markov equilibrium, prices depend on an aggregate state vector that includes $s_t$; thus $(r_t,\mu_{R_i,t},\sigma_{R_i,t})$ are functions of $(\ldots,s_t)$.

Absence of arbitrage opportunities implies the risk premium is given by
\begin{equation}
    \pi_{D,t} = \sigma_{R_D,t} \lambda_t', \qquad \qquad 
	\pi_{B,t} = \sigma_{R_B,t} \lambda_t',
\end{equation}
where $\lambda_t$ is the market price of risk.

Since $D_t+B_t=Y_t$, the combined claim replicates a claim to the aggregate endowment. Denoting the endowment-claim price by $P_{Y,t}$, we must have
\begin{equation}\label{eq:replication_prices}
    P_{Y,t}=P_{D,t}+P_{B,t}.
\end{equation}

\paragraph{Derivative claim.}
Let $\sigma_{R_F,t}$ be the $1\times 3$ vector of risk exposure of the derivative claim. 
We assume that the last component of $\sigma_{R_F,t}$ is equal to zero, $\sigma_{R_F,3,t} = 0$. 
As before, we assume that $\sigma_{R_F,t}$ is orthogonal to $\sigma_{R_Y,t}$, $\sigma_{R_F,t} \sigma_{R_Y,t}' = 0$, and we adopt the normalization $\|\sigma_{R_F,t} \| = \|\sigma_{R_Y,t} \|$.

\paragraph{Price of risk.} Let $\Sigma_t$ denote the matrix of risk exposures:
\begin{equation}
    \Sigma_t = \begin{pmatrix}
        \sigma_{R_B,t} \\
		\sigma_{R_D,t} \\
		 \sigma_{R_F,t}
    \end{pmatrix}.
\end{equation}
and $\pi_t \equiv [\pi_{B,t}, \pi_{D,t}, \pi_{F,t} ]'$.

The no-arbitrage condition can be written as
\begin{equation}
    \pi_t = \Sigma_t \lambda_t' \iff \lambda_t'  = \Sigma_t^{-1} \pi_t.
\end{equation}

\paragraph{Investors' problem.} 

An investor of type $j = 0,\ldots,J$ solves
\begin{equation}\label{eq:inv-problem}
    V_{j,t} = \max_{[C_j,\alpha_j,\theta_j]} 
    \mathbb{E}_t \left[ \int_t^\infty f_j(C_{j,s},V_{j,s})\, ds \right], \qquad s.t.
\end{equation}
% subject to the budget constraint  
\begin{equation}
dW_{j,t} = \big[(r_t + \pi_{B,t} \alpha_{j,B,t}+ \pi_{D,t} \alpha_{j,D,t} + \pi_{F,t} \theta_{j,t}) W_{j,t} - C_{j,t}\big] dt 
    + \Sigma_{j,t} W_{j,t} \, dZ_t,
\end{equation}
where 
\begin{equation}
    \Sigma_{j,t} = \alpha_{j,D,t} \sigma_{R_D,t} +\alpha_{j,B,t} \sigma_{R_B,t} + \theta_{j,t} \sigma_{R_F,t},
\end{equation}
with $W_{j,t} \geq 0$ and $(\alpha_{j,D,t}, \alpha_{j,B,t}, \theta_{j,t}) \in \Omega_{j,t}$. 

For passive investors ($j=0$),  
$\Omega_{0,t} = \{(\alpha_{0,D,t}, \alpha_{0,B,t}, \theta_{0,t}): \alpha_{0,D,t} = \overline{\alpha}_{p,t} \omega_{D,t}, \alpha_{0,B,t} = \overline{\alpha}_{p,t} \omega_{B,t}, \theta_{0,t} = 0\}$, where $\overline{\alpha}_{p,t}$ follows 
\begin{equation}
    d \overline{\alpha}_{p,t} = \theta_p (\overline{\alpha} - \overline{\alpha}_{p,t}) dt + \sigma_{p}  d Z_t,
\end{equation}	
and $\omega_{D,t} \equiv \frac{P_{D,t}}{P_{D,t}+P_{B,t}}$ and $\omega_{B,t} \equiv \frac{P_{B,t}}{P_{D,t}+P_{B,t}}$ are the market weights of the dividend and bond claims.
For active investors ($j=1,\ldots,J$),  
$\Omega_{j,t} = \{(\alpha_j, \theta_j): \| \Sigma_{j,t} \| \leq \overline{\sigma}\}$. 

\paragraph{Market clearing.} The  market clearing conditions are
\begin{equation}
	\sum_{j=0}^J \omega_j C_{j,t} = Y_t, 
	\qquad 
	\sum_{j=0}^J \omega_j Q_{j,D,t} = 1, 
	\qquad 
	\sum_{j=0}^J \omega_j Q_{j,B,t} = 1, 
	\qquad 
	\sum_{j=0}^J \omega_j B_{j,t} = 0,
	\qquad
	\sum_{j=0}^J \omega_j W_{j,t} \theta_{j,t} = 0,
\end{equation}
where $Q_{j,D,t} = \alpha_{j,D,t} W_{j,t} / P_{D,t}$ and $Q_{j,B,t} = \alpha_{j,B,t} W_{j,t} / P_{B,t}$ denote holdings of the dividend and bond claims. 

\subsection{Characterization of the equilibrium}

\paragraph{Value function and consumption-wealth ratio.} With homothetic preferences, investor $j$'s value function can be written as
\begin{equation}\nonumber
    V_{j,t} = \left( \frac{c_{j,t}}{\rho^\psi} \right)^{\frac{1-\gamma_j}{1-\psi}} 
    \frac{W_{j,t}^{1-\gamma_j}}{1-\gamma_j}, \qquad \text{where} \qquad \frac{dc_{j,t}}{c_{j,t}} = \mu_{c_j,t}\, dt + \sigma_{c_j,t}\, dZ_t,
\end{equation}
and $c_{j,t} = c_j(X_t)$ corresponds to the consumption-wealth ratio: $\frac{C_{j,t}}{W_{j,t}} = c_{j,t}$.

From the Hamilton-Jacobi-Bellman equation, the consumption-wealth ratio is
\begin{equation}\label{eq:HJB_complete}
    c_{j,t} = \psi \rho 
    + (1-\psi) \left[ r_t + \pi_{D,t} \alpha_{j,D,t} + \pi_{B,t} \alpha_{j,B,t} + \pi_{F,t} \theta_{j,t}
        - \frac{\gamma_j}{2} \|\Sigma_{j,t}\|^2 \right] 
    + \xi_{j,t},
\end{equation}
where 
\[
    \xi_{j,t} \equiv \mu_{c_j,t} 
    + (1-\gamma_j) \sigma_{c_j,t} \Sigma_{j,t}' 
    + \frac{\psi - \gamma_j}{1-\psi} \frac{\|\sigma_{c_j,t}\|^2}{2}.
\]

\paragraph{Optimal portfolio share and derivative exposure.} 
The active investor chooses $(\alpha_{j,D,t},\alpha_{j,B,t},\theta_{j,t})$ to maximize the right-hand side of \eqref{eq:HJB_complete} subject to $\|\Sigma_{j,t}\|\le\overline{\sigma}$. 
Denote the Lagrange multiplier on the leverage constraint by $\nu_{j,t}\ge 0$. The FOC with respect to $\alpha_{j,D,t}$ and $\alpha_{j,B,t}$ are respectively:
\begin{align}\label{eq:FOC_alpha_D_complete}
    (1-\psi)\Big[\pi_{D,t} - \gamma_j\,\Sigma_{j,t}\,\sigma_{R_D,t}'\Big] + (1-\gamma_j)\,\sigma_{c_j,t}\,\sigma_{R_D,t}' &= \nu_{j,t}\,\Sigma_{j,t}\,\sigma_{R_D,t}'\\
    (1-\psi)\Big[\pi_{B,t} - \gamma_j\,\Sigma_{j,t}\,\sigma_{R_B,t}'\Big] + (1-\gamma_j)\,\sigma_{c_j,t}\,\sigma_{R_B,t}' &= \nu_{j,t}\,\Sigma_{j,t}\,\sigma_{R_B,t}'.
\end{align}
The FOC with respect to $\theta_{j,t}$ is
\begin{equation}\label{eq:FOC_theta_complete}
    (1-\psi)\Big[\pi_{F,t} - \gamma_j\,\Sigma_{j,t}\,\sigma_{R_F,t}'\Big] + (1-\gamma_j)\,\sigma_{c_j,t}\,\sigma_{R_F,t}' = \nu_{j,t}\,\Sigma_{j,t}\,\sigma_{R_F,t}'.
\end{equation}

Using the fact that $\pi_{D,t} = \sigma_{R_D,t} \lambda_t'$, $\pi_{B,t} = \sigma_{R_B,t} \lambda_t'$, and $\pi_{F,t} = \sigma_{R_F,t} \lambda_t'$, we can rewrite the FOC as follows:
\begin{equation}
	\left[ \lambda_t - \lambda_{j,t}  \right] \sigma_{R_B,t}'  = 0, \qquad 
	\left[ \lambda_t - \lambda_{j,t}  \right] \sigma_{R_D,t}'  = 0, \qquad 
	\left[ \lambda_t - \lambda_{j,t}  \right] \sigma_{R_F,t}'  = 0,
\end{equation}
where
\begin{equation}
	\lambda_{j,t} \equiv \gamma_j \left[\Sigma_{j,t} + \frac{1 - \gamma_j^{-1}}{1-\psi} \sigma_{c_j,t}\right] + \frac{\nu_{j,t}}{1-\psi} \Sigma_{j,t}.
\end{equation}

The optimality conditions above imply that
\begin{equation}
    \lambda_{j,t} = \lambda_t \iff \lambda_{j,t}' = \Sigma_t^{-1} \pi_t.
\end{equation}

\subsection{Verifying the conjecture}

Let $(c_{j,t}^ *, \alpha_{j,t}^ *, \theta_{j,t}^ *)$ denote the solution to the investors' problem in the economy with the aggregate claim.
We conjecture that $c_{j,t} = c_{j,t}^*$, $\alpha_{j,D,t} = \omega_{D,t} \alpha_{j,t}^*$, $\alpha_{j,B,t} = \omega_{B,t} \alpha_{j,t}^*$, and $\theta_{j,t} = \theta_{j,t}^*$ for all $j$.
This implies that $\Sigma_{j,t} = (\Sigma_{j,t}^ *, 0)$, as investors are not exposed to the dividend share shock in the baseline model.
Similarly, we have $\sigma_{c_j,t} = (\sigma_{c_j,t}^*,0)$ and $\lambda_{j,t} = (\lambda_{j,t}^*,0)$.

Using the fact that $\pi_{Y,t} = \omega_{D,t} \pi_{D,t} + \omega_{B,t} \pi_{B,t}$ and 
$\sigma_{R_Y,t} = \omega_{D,t} \sigma_{R_D,t} + \omega_{B,t} \sigma_{R_B,t}$, we have that $c_{j,t} = c_{j,t}^*$ for all $j$ satisfies the HJB equation.
Given $\Sigma_{j,t} = (\Sigma_{j,t}^*,0)$, and the fact that the third entry of $\sigma_{c_j,t}$ is zero, we have that the first-order condition for the derivative is satisfied.


Using the fact that $\pi_{D,t} = \sigma_{R_D,t} \lambda_t'$ and $\pi_{B,t} = \sigma_{R_B,t} \lambda_t'$, and that the third entry of $\lambda_t$ is zero, we have that the first-order condition for the equity and debt claims are satisfied.

The market clearing conditions for the equity and debt claims can be written as
\begin{equation}
	\sum_{j=0}^J x_{j,t} \alpha_{j,D,t} = \omega_{D,t}, 
	\qquad 
	\sum_{j=0}^J x_{j,t} \alpha_{j,B,t} = \omega_{B,t}
\end{equation}

Using the fact that $\alpha_{0,D,t} = \overline{\alpha}_{p,t} \omega_{D,t}$, we obtain
\begin{equation}
	\left(x_{0,t}\overline{\alpha}_{p,t} + \sum_{j=1}^J x_{j,t} \alpha_{j,t}\right) \omega_{k,t} = \omega_{k,t} \Rightarrow x_{0,t}\overline{\alpha}_{p,t} + \sum_{j=1}^J x_{j,t} \alpha_{j,t} = 1,
\end{equation}
for $k = \{D, B\}$. Hence, the market clearing conditions for the equity and debt claims are satisfied.

Our conjecture implies that the remaining market clearing conditions are also satisfied.
This completes the verification of the conjecture.

\subsection{Pricing the dividend and risky debt claims}

Let $\hat X_t = (X_t,s_t)$ denote the vector of aggregate state variables in the model with the dividend and risky-debt claims, where $X_t$ is the vector of aggregate state variables in the baseline model.

From the baseline model (which only includes the aggregate endowment claim), the stochastic discount factor satisfies
\begin{equation}
\frac{d\eta_t}{\eta_t} = -r(X_t)dt - \lambda(X_t)dZ_t,
\end{equation}
where $r(X_t)$ is the short rate and $\lambda(X_t)$ is the market price of risk. Importantly, these objects depend only on $X_t$ and are therefore independent of the dividend share $s_t$. This property allows us to solve first for the equilibrium dynamics of $X_t$ and the SDF in the baseline economy and then price individual claims given these dynamics.

The pricing condition for each asset $k \in \{D,B\}$ is
\begin{equation}
\mu_{R_k,t} - r_t = \sigma_{R_k,t}\lambda_t',
\end{equation}
where diffusion terms are written as row vectors.

\paragraph{Valuation ratios.}

Let $p_{k,t} \equiv P_{k,t}/Y_t$ denote the valuation ratio of claim $k$, where $Y_t$ follows
\begin{equation}
\frac{dY_t}{Y_t} = \mu dt + \sigma dZ_t,
\end{equation}
with constant $\mu$ and $\sigma$ (with $\sigma$ a $1\times d$ row vector). Since the valuation ratio may depend on both the aggregate state and the dividend share, we write
\begin{equation}
p_{k,t} = p_k(X_t,s_t).
\end{equation}

The state variables evolve according to
\begin{align}
dX_t &= \mu_X(X_t)dt + \sigma_X(X_t)dZ_t, \\
ds_t &= \mu_s(s_t)dt + \sigma_s(s_t)dZ_t,
\end{align}
where $\sigma_X(X_t)$ is an $m \times d$ diffusion matrix and $\sigma_s(s_t)$ is a $1\times d$ row vector. In our specification
\begin{equation}
\mu_s(s_t) = \kappa_s(\bar s - s_t),
\qquad
\sigma_s(s_t) = s_t(1-s_t)\bar\sigma_s,
\end{equation}
with constant $\bar\sigma_s$.

Using the convention that derivatives are row vectors, define
\[
p_X \equiv \frac{\partial p_k}{\partial X}\in\mathbb{R}^{1\times m},
\qquad
p_s \equiv \frac{\partial p_k}{\partial s}\in\mathbb{R},
\]
and let $p_{XX}\in\mathbb{R}^{m\times m}$, $p_{Xs}\in\mathbb{R}^{1\times m}$, and $p_{ss}\in\mathbb{R}$ denote the corresponding second derivatives.

Applying It\^{o}'s lemma to $p_{k,t}=p_k(X_t,s_t)$ yields
\begin{equation}
\frac{dp_{k,t}}{p_{k,t}} = \mu_{p_k,t}dt + \sigma_{p_k,t}dZ_t,
\end{equation}
where the diffusion term is
\begin{equation}
\sigma_{p_k,t}
=
\frac{1}{p_{k,t}}
\left(
p_X\sigma_X + p_s\sigma_s
\right)
\end{equation}
and the drift term is
\begin{align}
\mu_{p_k,t}
=
\frac{1}{p_{k,t}}
\Bigg[
&p_X\mu_X + p_s\mu_s \notag\\
&+\frac12
\Big(
\mathrm{tr}(\sigma_X\sigma_X' p_{XX})
+2\,p_{Xs}(\sigma_X\sigma_s')
+(\sigma_s\sigma_s')p_{ss}
\Big)
\Bigg].
\end{align}
The cross term $p_{Xs}(\sigma_X\sigma_s')$ reflects the instantaneous covariance $d\langle X,s\rangle_t=\sigma_X\sigma_s' dt$; it disappears if $\sigma_X\sigma_s'=0$.

\paragraph{Returns.}

Since $P_{k,t}=p_{k,t}Y_t$, the return volatility is
\begin{equation}
\sigma_{R_k,t} = \sigma + \sigma_{p_k,t},
\end{equation}
and the expected return is
\begin{equation}
\mu_{R_k,t}
=
\frac{s_{k,t}}{p_{k,t}}
+
\mu
+
\mu_{p_k,t}
+
\langle \sigma, \sigma_{p_k,t} \rangle,
\qquad
\langle a,b\rangle \equiv a b'.
\end{equation}
Here $s_{D,t}=s_t$ and $s_{B,t}=1-s_t$.

Substituting these expressions into the pricing condition
$\mu_{R_k,t}-r_t=\sigma_{R_k,t}\lambda_t'$ yields the partial differential equation for the valuation ratio $p_k(X,s)$:
\begin{align}
0=&\;
p_X\,\mu_X(X)+p_s\,\mu_s(s) \notag\\
&+
\frac12
\Big[
\mathrm{tr}(\sigma_X(X)\sigma_X(X)'p_{XX})
+
2\,p_{Xs}(\sigma_X(X)\sigma_s(s)')
+
(\sigma_s(s)\sigma_s(s)')p_{ss}
\Big] \notag\\
&-(r(X)-\mu)p_k(X,s)
+
s_k(s) \notag\\
&+
p_k(X,s)\,\sigma\,\sigma_{p_k}(X,s)' \notag\\
&-
p_k(X,s)\,\sigma\,\lambda(X)' \notag\\
&-
p_k(X,s)\,\sigma_{p_k}(X,s)\lambda(X)' .
\end{align}

% \paragraph{Risk-neutral generator form.}

% Define the risk-neutral drifts
% \[
% \tilde\mu_X(X) \equiv \mu_X(X) - \sigma_X(X)\lambda(X)',
% \qquad
% \tilde\mu_s(X,s) \equiv \mu_s(s) - \sigma_s(s)\lambda(X)'.
% \]
% Let the corresponding generator acting on $p_k(X,s)$ be
% \begin{align}
% \tilde{\mathcal L}p_k
% \equiv&\;
% p_X\,\tilde\mu_X(X) + p_s\,\tilde\mu_s(X,s) \notag\\
% &+\frac12\Big[
% \mathrm{tr}(\sigma_X(X)\sigma_X(X)' p_{XX})
% +2\,p_{Xs}(\sigma_X(X)\sigma_s(s)')
% +(\sigma_s(s)\sigma_s(s)')p_{ss}
% \Big].
% \end{align}
% Then the valuation ratio $p_k(X,s)$ solves
% \begin{equation}
% \tilde{\mathcal L}p_k(X,s)
% -(r(X)-\mu+\sigma\lambda(X)')\,p_k(X,s)
% +s_k(s)
% +p_k(X,s)\,(\sigma-\lambda(X))\,\sigma_{p_k}(X,s)'
% =0,
% \end{equation}
% where
% \[
% \sigma_{p_k}(X,s)=\frac{1}{p_k(X,s)}\Big(p_X\sigma_X(X)+p_s\sigma_s(s)\Big).
% \]


\newpage
\singlespacing
%\bibliographystyle{jf}
\phantomsection\addcontentsline{toc}{section}{\refname}
\clearpage
%{\small
\bibliography{../references.bib}
\end{document}